% ---
% titre: Aire d'une façade de maison (figure complexe)
% theme: GM
% chapitre: Lignes et surfaces
% sous_chapitre: Lignes et surfaces
% niveau: [9e]
% type: EVAL
% tags: [c-perimetre, c-aire, c-figure-complexe, p-question-TS]
% difficulte: 2
% corrige: true
% ---

\question[10]
Alexandre doit repeindre la façade de sa maison, représentée ci-dessous. Les deux fenêtres et la porte ne doivent pas être peintes.

\vspace{10pt}
\begin{center}
\begin{tikzpicture}[scale=1, dim style/.style={black,dashed}, dim fence style/.style={draw=none}]
    \tkzDefPoints{0/0/A, 4/0/B, 4/3/C, 0/3/D, 2/4.5/E, 0.7/1.6/F, 1.5/1.6/G, 1.5/2.4/H, 0.7/2.4/I,  2.5/1.6/J, 3.3/1.6/K, 3.3/2.4/L, 2.5/2.4/M, 1.6/0/N, 2.6/0/O, 2.6/1.3/P, 1.6/1.3/Q, 2/3/R}
    % corps de la maison + toit
    \tkzDrawSegments(A,B B,C D,A)
    \tkzDrawSegments(C,E D,E)
    % fenêtres
    \tkzDrawPolygon(F,G,H,I)
    \tkzDrawPolygon(J,K,L,M)
    % porte
    \tkzDrawPolygon(N,O,P,Q)

    \tkzLabelPoint[below left](A){$A$}
    \tkzLabelPoint[below right](B){$B$}
    \tkzLabelPoint[right](C){$C$}
    \tkzLabelPoint[left](D){$D$}
    \tkzLabelPoint[above](E){$E$}

%    \tkzDrawSegments[dashed](E,R)
    \tkzDrawSegment[draw=none, dim={{1,5\,m},-80pt,midway}, dim style/.append style={black,dashed},dim fence style/.append style={draw=none}](E,R)
    \tkzDrawSegment[draw=none, dim={{4\,m},-20pt,midway}, dim style/.append style={black,dashed},dim fence style/.append style={draw=none}](A,B)
    \tkzDrawSegment[draw=none, dim={{3\,m},-20pt,midway}, dim style/.append style={black,dashed},dim fence style/.append style={draw=none}](B,C)
\end{tikzpicture}
\end{center}

\vspace{5pt}
Les deux fenêtres sont des carrés de 0,8\,m de côté. La porte mesure 1\,m sur 2\,m.

\begin{parts}
\part[10] Calcule l'aire de la façade à peindre.

\vspace{10pt}{\footnotesize\raggedleft Espace pour ta démarche et tes calculs\par}
\begin{solutionorgrid}[220pt]
$A_{corps} = 4\cdot 3 = 12\,m^2$ \hfill \textbf{(2pt)}

\vspace{5pt}
$A_{toit} = \dfrac{4\cdot 1,5}{2} = 3\,m^2$ \hfill \textbf{(2pt)}

\vspace{5pt}
$A_{fenêtres} = 2\cdot 0,8^2 = 2\cdot 0,64 = 1,28\,m^2$ \hfill \textbf{(2pt)}

\vspace{5pt}
$A_{porte} = 1\cdot 2 = 2\,m^2$ \hfill \textbf{(1,5pt)}

\vspace{5pt}
$A_{peinte} = A_{corps}+A_{toit}-A_{fenêtres}-A_{porte} = 12+3-1,28-2 = 11,72\,m^2$ \hfill \textbf{(2pt)}
\end{solutionorgrid}

\begin{tcolorbox}[enhanced jigsaw, interior hidden, colframe=box, parbox=false]%
Réponse: 
\sol{\vspace{20pt}}{L'aire de la façade à peindre vaut 11,72\,m$^2$. \textbf{(0,5pt)}}
\end{tcolorbox}

\newpage
\part[5]
Pour Noël, Alexandre veut décorer sa maison et installe une guirlande lumineuse tout autour des deux fenêtres et de la porte.

Quelle est la longueur de guirlande nécessaire ?

\vspace{10pt}{\footnotesize\raggedleft Espace pour ta démarche et tes calculs\par}
\begin{solutionorgrid}[180pt]
La guirlande fait le tour de chaque ouverture séparément : il faut donc calculer le périmètre de chaque fenêtre et de la porte. \hfill \textbf{(1pt)}

\vspace{5pt}
$P_{fenêtre} = 4\cdot 0,8 = 3,2\,m$ \hfill \textbf{(1pt)}

\vspace{5pt}
Pour les deux fenêtres : $2\cdot 3,2 = 6,4\,m$ \hfill \textbf{(1pt)}

\vspace{5pt}
$P_{porte} = 2\cdot(1+2) = 6\,m$ \hfill \textbf{(1pt)}

\vspace{5pt}
$P_{total} = 6,4 + 6 = 12,4\,m$ \hfill \textbf{(1pt)}
\end{solutionorgrid}

\begin{tcolorbox}[enhanced jigsaw, interior hidden, colframe=box, parbox=false]%
Réponse: 
\sol{\vspace{20pt}}{Alexandre a besoin de 12,4\,m de guirlande.}
\end{tcolorbox}
\end{parts}